\renewcommand{\thetheorem}{15.a.\arabic{theorem}}
\setcounter{theorem}{0}

\setcounter{chapter}{14} % Sets chapter counter so the current chapter is 15
\chapter{Linear Maps}

% ---------------------------------------------------------------------------
% Provenance of the prose in this file.
%   % Extractor: <model>   the block below was transcribed from Prof. Biran's
%                          handwritten notes by that model
%   % Creator:   <model>   the block below was written by that model and has no
%                          counterpart in the notes
% A tag holds until the next tag.
% ---------------------------------------------------------------------------

% Extractor: Gemini 3.1 Pro
\section{Introduction to Linear Maps}
\label{sec:intro_linear_maps}

Linear maps, also known as linear transformations, are homomorphisms of vector spaces. They are the primary objects of study in linear algebra, as they preserve the underlying algebraic structure of vector spaces.

% def:15.a.1
\begin{definition}[Linear Map]
\label{def:linear_map}
Let $V$ and $W$ be vector spaces over a field $K$. A map $T: V \to W$ is called a \newterm{linear map} if it satisfies the following two conditions:
\begin{enumerate}
    \item For all $v_1, v_2 \in V$, $T(v_1 + v_2) = T(v_1) + T(v_2)$.
    \item For all $\alpha \in K$ and $v \in V$, $T(\alpha \cdot v) = \alpha \cdot T(v)$.
\end{enumerate}
\end{definition}

\begin{remark}[Notation and Terminology]
\label{rem:linear_map_notation}
In many contexts, we omit the parentheses and write $Tv$ instead of $T(v)$. The conditions outlined in \cref{def:linear_map} essentially state that $T: V \to W$ is a map that respects the operations defined on the vector spaces; specifically, the addition and scalar multiplication in $V$ are mapped perfectly to the corresponding operations in $W$.

We denote the set of all linear maps from $V$ to $W$ by $\Hom_K(V, W)$, or simply $\Hom(V, W)$. While some authors use alternative notations such as $\operatorname{Lin}(V, W)$ or $\hom_K(V, W)$, $\Hom_K(V, W)$ remains the standard academic convention.
\end{remark}

\begin{exercise}[Linearity Criterion]
\label{exer:linearity_criterion}
Prove that a map $T: V \to W$ is linear if and only if for all $\alpha_1, \alpha_2 \in K$ and $v_1, v_2 \in V$:
\[
T(\alpha_1 v_1 + \alpha_2 v_2) = \alpha_1 T(v_1) + \alpha_2 T(v_2).
\]
\end{exercise}

\begin{exercise}[Extending a Linear Map by Zero Outside a Subspace]
\label{exc:extending_linear_map_by_zero}
Let $V$ and $W$ be vector spaces over a field $K$. Suppose that $U \subsetneq V$ is a proper linear subspace, and let $S \in \Hom_K(U, W)$ be a non-trivial linear map (i.e., $S(u) \ne 0_W$ for some $u \in U$). Define the extended map $T \colon V \to W$ by
\[
T(v) :=
\begin{cases}
S(v) & \text{if } v \in U, \\
0_W & \text{if } v \in V \setminus U.
\end{cases}
\]
Is $T$ a linear map? Justify your answer.
\exinfo{This exercise is Problem 2 of Problem Sheet 10, Linear Algebra I, Autumn Semester 2025.}
\end{exercise}

\begin{exercise}[Additive Maps over the Rational Field \texorpdfstring{$\mathbb{Q}$}{Q}]
\label{exc:additive_maps_over_rationals}
Let $V, W$ be vector spaces over $\mathbb{Q}$. A map $f \colon V \to W$ is called \newterm{additive} if
\[
f(x + y) = f(x) + f(y) \qquad \text{for all } x, y \in V.
\]
Show that every additive map is automatically $\mathbb{Q}$-linear, so that
\[
\Hom_{\mathbb{Q}}(V, W) = \{f \colon V \to W \mid f \text{ is additive}\}.
\]
\exinfo{This exercise is Problem 7 of Problem Sheet 10, Linear Algebra I, Autumn Semester 2025.}
\end{exercise}

\begin{example}[Fundamental Examples of Linear Maps]
\label{ex:fundamental_linear_maps}
\begin{enumerate}
    \item The identity map $\id_V: V \to V$, defined by $\id_V(v) := v$ for all $v \in V$.
    \item The zero map $0: V \to W$, defined by $0(v) := 0_W$ for all $v \in V$. 
    
    Clearly, these initial examples are elementary, yet they serve as vital benchmarks.
    
    \item The derivative map $D: K[x] \to K[x]$. For any polynomial $p(x) = \sum_{k=0}^n a_k x^k$, we define its image as:
    \[
    D(p(x)) := \sum_{k=1}^n k a_k x^{k-1}.
    \]
    This map is linear, as the derivative of a linear combination of polynomials is simply the same linear combination of their derivatives.
    
    \item Consider the vector space $C([a, b])$ consisting of continuous functions $f: [a, b] \to \mathbb{R}$. The definite integral defines a linear map $T: C([a, b]) \to \mathbb{R}$, given by:
    \[
    T(f) := \int_a^b f(x) \, dx.
    \]
    
    \item The shift map $S: K^\infty \to K^\infty$, defined on infinite sequences as:
    \[
    S((a_1, a_2, a_3, \dots)) := (a_2, a_3, a_4, \dots).
    \]
\end{enumerate}
\end{example}

\begin{example}[Matrix Transformation]
\label{ex:matrix_transformation}
This is a central example in finite-dimensional linear algebra. Let $A \in \M_{m \times n}(K)$ be a matrix. We define a map $T_A: K^n \to K^m$ via matrix-vector multiplication:
\[
T_A \left(\begin{smallmatrix} x_1 \\ \vdots \\ x_n \end{smallmatrix}\right) := A \cdot \left(\begin{smallmatrix} x_1 \\ \vdots \\ x_n \end{smallmatrix}\right).
\]
Recall from matrix arithmetic that $A \cdot (\alpha \cdot x) = \alpha \cdot (A \cdot x)$ and $A \cdot (x + y) = A \cdot x + A \cdot y$. Consequently, $T_A$ is a linear map.
\end{example}

\begin{exercise}[Linearity of Integral and Evaluation Map on Polynomials]
\label{exc:polynomial_integral_eval_linear_map}
Let $b, c \in \mathbb{R}$ and let $\mathbb{R}[x]$ denote the space of polynomial functions in one variable with real coefficients. Define the map $T \colon \mathbb{R}[x] \to \mathbb{R}^2$ by
\[
T(p) := \left( 3p(4) + 5p^3(6) + b p(1)p(2), \; \int_{-1}^2 x^2 p(x) \, dx + c (p(0))^2 \right)\transp.
\]
Show that $T$ is linear if and only if $b = c = 0$.
\exinfo{This exercise is Problem 2 of Problem Sheet 9, Linear Algebra I, Autumn Semester 2025.}
\end{exercise}

\begin{exercise}[Identifying Linear Maps]
\label{exc:identifying_linear_maps_mc}
Which of the following maps are linear?
\begin{enumerate}
    \item $f \colon \mathbb{R}^2 \to \mathbb{R}^3, \quad (x, y) \mapsto (x + y, 2x, 0)$;
    \item $f \colon \mathbb{R}^2 \to \mathbb{R}^3, \quad (x, y) \mapsto (x + y, 2xy, 0)$;
    \item $f \colon K^3 \to K^2, \quad (x, y, z) \mapsto (\alpha x + \beta y + \gamma z, \; \delta x + \varepsilon y + \eta z)$ for fixed constants $\alpha, \beta, \gamma, \delta, \varepsilon, \eta \in K$.
\end{enumerate}
\exinfo{This exercise is Multiple Choice Question 1 of Problem Sheet 9, Linear Algebra I, Autumn Semester 2025.}
\end{exercise}

\begin{exercise}[Matrix Representation of a Linear Map]
\label{exc:matrix_map_mc}
Consider the matrix $A = \begin{pmatrix} 2 & 1 \\ 1 & 0 \\ 0 & 2 \end{pmatrix}$. Which of the following linear maps is given by $x \mapsto A \cdot x$?
\begin{enumerate}
    \item $f \colon \mathbb{R}^2 \to \mathbb{R}^3, \quad f(x, y) = (2x + y, x, 2y)\transp$;
    \item $f \colon \mathbb{R}^3 \to \mathbb{R}^2, \quad f(x, y, z) = (2x + y, x + 2z)\transp$.
\end{enumerate}
\exinfo{This exercise is Multiple Choice Question 2 of Problem Sheet 9, Linear Algebra I, Autumn Semester 2025.}
\end{exercise}

% prop:15.a.2
\begin{proposition}[Basic Properties of Linear Maps]
\label{prop:properties_linear_maps}
Let $V$ and $W$ be vector spaces over a field $K$, and let $T: V \to W$ be a linear map. Then, the following properties hold:
\begin{enumerate}
    \item For all $n \in \mathbb{N}$, $\alpha_1, \dots, \alpha_n \in K$, and $v_1, \dots, v_n \in V$, we have:
    \[
    T\left( \sum_{i=1}^n \alpha_i v_i \right) = \sum_{i=1}^n \alpha_i T(v_i).
    \]
    \item $T(0_V) = 0_W$.
\end{enumerate}
\end{proposition}

\begin{proof}
\begin{enumerate}
    \item The first statement follows directly from the definition of linearity, \cref{def:linear_map}, by applying induction on the number of vectors $n$.
    \item To prove the second statement, we evaluate the image of the zero vector:
    \[
    0_W = T(0_V) - T(0_V) = T(0_V - 0_V) = T(0_V).
    \]
    Therefore, any linear map must map the zero vector of the domain directly to the zero vector of the codomain.
\end{enumerate}
\end{proof}

% thm:15.a.3
\begin{theorem}[Existence and Uniqueness of Linear Maps on a Basis]
\label{thm:existence_uniqueness_linear_map}
Let $V$ and $W$ be vector spaces over $K$ with $V$ finite-dimensional, and let $\mathcal{B} = (v_1, \dots, v_n)$ be a basis for $V$. Given any arbitrary sequence of vectors $w_1, \dots, w_n \in W$, there exists a unique linear map $T: V \to W$ such that:
\[
T(v_i) = w_i, \quad \forall \, 1 \le i \le n.
\]
\end{theorem}

\begin{proof}
\textbf{Existence:} Let $v \in V$ be an arbitrary vector. Since $\mathcal{B}$ is a basis, there exist unique scalars $a_1, \dots, a_n \in K$ such that $v = \sum_{i=1}^n a_i v_i$. We define the map $T$ by setting:
\begin{equation}
\label{eq:T_def_existence}
T(v) := \sum_{i=1}^n a_i w_i.
\end{equation}
Because the coordinates $a_i$ are uniquely determined for each $v$, the map $T$ is well-defined.

% Extractor: Opus 5
We claim that $T$ is linear. Indeed, let $v, v' \in V$. There exist $a_1, \dots, a_n \in K$ with $v = a_1 v_1 + \dots + a_n v_n$ and $a'_1, \dots, a'_n \in K$ with $v' = a'_1 v_1 + \dots + a'_n v_n$, so that
\[
v + v' = (a_1 + a'_1) v_1 + \dots + (a_n + a'_n) v_n.
\]
By our definition of $T$, this gives
\[
T(v + v') = (a_1 + a'_1) w_1 + \dots + (a_n + a'_n) w_n = (a_1 w_1 + \dots + a_n w_n) + (a'_1 w_1 + \dots + a'_n w_n) = T(v) + T(v').
\]
Also, if $\alpha \in K$, then $\alpha v = \alpha a_1 v_1 + \dots + \alpha a_n v_n$, and therefore,
\[
T(\alpha v) = \alpha a_1 w_1 + \dots + \alpha a_n w_n = \alpha (a_1 w_1 + \dots + a_n w_n) = \alpha T(v).
\]
This proves that $T: V \to W$ is a linear map.

We claim next that $T(v_i) = w_i$ for all $1 \le i \le n$. Indeed, writing $v_i = 0 \cdot v_1 + \dots + 1 \cdot v_i + \dots + 0 \cdot v_n$ exhibits the coordinates of $v_i$, whence
\[
T(v_i) = 0 \cdot w_1 + \dots + 1 \cdot w_i + \dots + 0 \cdot w_n = w_i.
\]
This concludes the proof of the existence of $T$.

% Extractor: Gemini 3.1 Pro

\textbf{Uniqueness:} Suppose there exists another linear map $S: V \to W$ such that $S(v_i) = w_i$ for all $1 \le i \le n$. For any $v = \sum_{i=1}^n a_i v_i \in V$, the linearity of $S$ dictates:
\[
S(v) = S\left( \sum_{i=1}^n a_i v_i \right) = \sum_{i=1}^n a_i S(v_i) = \sum_{i=1}^n a_i w_i = T(v).
\]
Since $S(v) = T(v)$ for all $v \in V$, the maps are functionally identical, establishing that $S = T$.
\end{proof}

\stepcounter{theorem}
% ex:15.a.4
\begin{example}[Linear Maps on Coordinate Spaces]
\label{ex:linear_map_std_basis}
Consider the coordinate space $K^n$ equipped with the standard basis $\mathcal{E} = (e_1, \dots, e_n)$. Let $w_1, \dots, w_n \in K^m$ be any vectors. \Cref{thm:existence_uniqueness_linear_map} guarantees a unique linear map $T: K^n \to K^m$ such that $T(e_i) = w_i$ for all $1 \le i \le n$.

\begin{question}
Can we provide an explicit matrix description for this map $T$?
\end{question}

\begin{answer}
Let $A \in \M_{m \times n}(K)$ be the matrix whose $j$-th column is precisely the vector $w_j$:
\[
A = 
\begin{pmatrix}
| & & | \\
w_1 & \dots & w_n \\
| & & |
\end{pmatrix}.
\]
The linear transformation $T_A: K^n \to K^m$ defined by $v \mapsto A \cdot v$ satisfies $T_A(e_j) = w_j$ for all $j$.
% Extractor: Opus 5
Indeed, multiplying $A$ by $e_i$ picks out the $i$-th column of $A$:
\[
T_A(e_i) = \begin{pmatrix} | & & | & & | \\ w_1 & \dots & w_i & \dots & w_n \\ | & & | & & | \end{pmatrix} \cdot \begin{pmatrix} 0 \\ \vdots \\ 1 \\ \vdots \\ 0 \end{pmatrix} = \begin{pmatrix} | \\ w_i \\ | \end{pmatrix},
\]
where the $1$ on the right sits in entry $i$.
% Extractor: Gemini 3.1 Pro
Due to the absolute uniqueness property established in \cref{thm:existence_uniqueness_linear_map}, we immediately conclude that $T = T_A$.
\end{answer}
\end{example}

% lem:15.a.5
\begin{lemma}[Matrix Representation of Linear Maps]
\label{lem:all_linear_maps_are_matrices}
For every linear transformation $T: K^n \to K^m$, there exists a unique matrix $A \in \M_{m \times n}(K)$ such that $T = T_A$.
\end{lemma}

\begin{proof}
Let $w_i := T(e_i)$ for each $i \in \{1, \dots, n\}$. As demonstrated in the preceding example, the matrix $A$ formed by using these specific vectors as columns defines a linear map $T_A$ that perfectly agrees with $T$ on the standard basis $\mathcal{E}$. Because $\mathcal{E}$ is a basis for $K^n$, \cref{thm:existence_uniqueness_linear_map} strictly implies that $T$ and $T_A$ must be the exact same map.

% Creator: Opus 5
The matrix is moreover unique, which the notes do not claim but which costs one line: if $B \in \M_{m \times n}(K)$ also satisfies $T = T_B$, then the $i$-th column of $B$ equals $B \cdot e_i = T_B(e_i) = T(e_i) = w_i$ for every $i$, so $B$ and $A$ have the same columns and $B = A$.
\end{proof}
% Source: 15.linear maps.b ("Kernel and Image"), which numbers from 15.b.1.
\renewcommand{\thetheorem}{15.b.\arabic{theorem}}
\setcounter{theorem}{0}
\section{Kernel, Image, and the Rank-Nullity Theorem}
\label{sec:kernel_image_rank}

In this section, we explore two fundamental subspaces associated with every linear map: the kernel and the image. These subspaces provide deep insight into the structure of the transformation and ultimately lead us to one of the most celebrated results in linear algebra, the Rank-Nullity Theorem.

% prop:15.b.1
\begin{proposition}[Kernel and Image]
\label{prop:kernel_image}
Let $V$ and $W$ be vector spaces over a field $K$, and let $T: V \to W$ be a linear map.
\begin{enumerate}[label=\textbf{(\alph*)}]
    \item The set $\ker(T) := \{v \in V \mid T(v) = 0_W\} \subseteq V$ is a linear subspace of $V$. It is called the \newterm{kernel} of $T$.
    \item The set $\Img(T) := \{T(v) \mid v \in V\} \subseteq W$ is a linear subspace of $W$. It is called the \newterm{image} of $T$.
\end{enumerate}
\end{proposition}

\begin{remark}[Set-Theoretic Perspective]
\label{rem:set_theoretic_notation}
Using standard notation from set theory, we can concisely write the kernel as the fiber over the zero vector, $\ker(T) = T^{-1}(\{0_W\})$, and the image as the range of the entire domain, $\Img(T) = T(V)$. (Note: The image is sometimes denoted as $\operatorname{Image}(T)$).
\end{remark}

\begin{proof}[Proof of \Cref{prop:kernel_image}]
\begin{enumerate}[label=\textbf{(\alph*)}]
    \item To show that $\ker(T) \subseteq V$ is a subspace, let $a, b \in K$, and let $v_1, v_2 \in \ker(T)$. By the linearity of $T$, we have:
    \[
    T(a v_1 + b v_2) = a T(v_1) + b T(v_2) = a \cdot 0_W + b \cdot 0_W = 0_W.
    \]
    This implies that $a v_1 + b v_2 \in \ker(T)$. Furthermore, since $T(0_V) = 0_W$, the zero vector $0_V \in \ker(T)$. Thus, the kernel is a valid subspace.
    
    \item To show that $\Img(T) \subseteq W$ is a subspace, let $a, b \in K$, and let $w_1, w_2 \in \Img(T)$. By definition, there exist vectors $v_1, v_2 \in V$ such that $T(v_1) = w_1$ and $T(v_2) = w_2$. Therefore, we can write:
    \[
    a w_1 + b w_2 = a T(v_1) + b T(v_2) = T(a v_1 + b v_2).
    \]
    Since $a v_1 + b v_2 \in V$, it follows immediately that $a w_1 + b w_2 \in \Img(T)$. Moreover, since $T(0_V) = 0_W$, we have $0_W \in \Img(T)$. Thus, the image is a subspace.
\end{enumerate}
\end{proof}

% Extractor: Opus 5
\begin{example}[The Kernel of a Matrix Transformation]
\label{ex:kernel_of_T_A}
Let $A \in \M_{m \times n}(K)$ and consider the linear map $T_A: K^n \to K^m$, recalling that $T_A(x) = A \cdot x$ for all $x \in K^n$. Then
\[
\ker(T_A) = \{x \in K^n \mid A \cdot x = 0\},
\]
which is precisely the solution set of the homogeneous linear system with coefficient matrix $A$. The kernel is therefore not a new object at all: it is the solution set of a homogeneous system, seen from the point of view of linear maps.
\end{example}

% prop:15.b.2
\begin{proposition}[Injectivity, Surjectivity, Kernel and Image]
\label{prop:injectivity_kernel}
Let $T: V \to W$ be a linear map. Then:
\begin{enumerate}[label=\textbf{(\alph*)}]
    \item $T$ is injective if and only if $\ker(T) = \{0_V\}$.
    \item $T$ is surjective if and only if $\Img(T) = W$.
\end{enumerate}
Recall that a map is called \newterm{bijective} if it is both injective and surjective.
\end{proposition}

% Extractor: Gemini 3.1 Pro
\begin{proof}
% Extractor: Opus 5
Statement \textbf{(b)} is a restatement of the definition of surjectivity, since $\Img(T) = T(V)$, so there is nothing to prove. We turn to \textbf{(a)}.

% Extractor: Gemini 3.1 Pro
First, suppose $T$ is injective. If $v \in \ker(T)$, then $T(v) = 0_W$. Since we also know that $T(0_V) = 0_W$, the injectivity of $T$ forces $v = 0_V$. Thus, $\ker(T) = \{0_V\}$.

Conversely, suppose that $\ker(T) = \{0_V\}$. Let $v_1, v_2 \in V$ such that $T(v_1) = T(v_2)$. By linearity, we obtain:
\[
T(v_1 - v_2) = T(v_1) - T(v_2) = 0_W.
\]
This implies that $v_1 - v_2 \in \ker(T)$. By our assumption, the only element in the kernel is the zero vector, and therefore, $v_1 - v_2 = 0_V$, which yields $v_1 = v_2$. Consequently, $T$ is injective.
\end{proof}

% Extractor: Opus 5
\begin{exercise}[Images and Preimages of Subspaces]
\label{exc:images_preimages_subspaces}
Let $T: V \to W$ be a linear map. Show that:
\begin{enumerate}[label=\textbf{(\alph*)}]
    \item if $V' \subseteq V$ is a linear subspace, then $T(V') \subseteq W$ is a linear subspace;
    \item if $W' \subseteq W$ is a linear subspace, then $T^{-1}(W') \subseteq V$ is a linear subspace, where $T^{-1}(W') := \{v \in V \mid T(v) \in W'\}$;
    \item both statements are generalisations of \cref{prop:kernel_image}. Explain why.
\end{enumerate}
\end{exercise}

\begin{exercise}[Preimage of Subspaces and Dimension Formula]
\label{exc:preimage_subspace_dimension_formula}
Let $f \colon V \to W$ be a linear map of vector spaces over $K$.
\begin{enumerate}
    \item Show that for every linear subspace $W' \subseteq W$, the preimage
    \[
    f^{-1}(W') := \{v \in V \mid f(v) \in W'\}
    \]
    is a linear subspace of $V$.
    \item If $V$ is finite-dimensional, prove the dimension formula:
    \[
    \dim f^{-1}(W') = \dim \ker(f) + \dim\big(\Img(f) \cap W'\big).
    \]
\end{enumerate}
\exinfo{This exercise is Problem 5 of Problem Sheet 10, Linear Algebra I, Autumn Semester 2025.}
\end{exercise}

\subsection{Isomorphisms}

Recall from set theory that a map $f: X \to Y$ is called bijective if it is both injective and surjective, and that in this case there exists a unique inverse map $g: Y \to X$, denoted by $f^{-1}$, meaning that
\[
g \circ f = \id_X \quad \text{and} \quad f \circ g = \id_Y.
\]
Conversely, if $f: X \to Y$ is a map for which some $g: Y \to X$ satisfies $g \circ f = \id_X$ and $f \circ g = \id_Y$, then $f$ is bijective. What is the right analogue of this for linear maps?

% def:15.b.3
\begin{definition}[Isomorphism]
\label{def:isomorphism}
A linear map $T: V \to W$ is called an \newterm{isomorphism} if there exists a \emph{linear} map $S: W \to V$ such that
\[
S \circ T = \id_V \quad \text{and} \quad T \circ S = \id_W.
\]
% Extractor: Gemini 3.1 Pro
Two vector spaces $V$ and $W$ are defined to be \qt{isomorphic}, denoted as $V \cong W$, if there exists an isomorphism between them.

\begin{center}
\begin{tikzcd}[row sep=large, column sep=huge]
V \arrow[r, "T", shift left=1.2ex] & W \arrow[l, "S", shift left=1.2ex]
\end{tikzcd}
\end{center}

\end{definition}

% Extractor: Opus 5
\begin{remark}
\label{rem:iso_is_bijective}
An isomorphism $T: V \to W$ is obviously bijective: the linear map $S$ of \cref{def:isomorphism} is in particular an inverse map of $T$ in the set-theoretic sense recalled above. The converse is the substantial question, and it is not a matter of definition, because the definition of an isomorphism demands that the inverse be \emph{linear}: is every bijective linear map an isomorphism?
\end{remark}

% lem:15.b.4
\begin{lemma}[Bijective Linear Maps are Isomorphisms]
\label{lem:bijective_is_iso}
Let $T: V \to W$ be a linear map which is bijective. Then the inverse map $T^{-1}: W \to V$ is linear. In other words, a bijective linear map is an isomorphism.
\end{lemma}

\begin{proof}
Let $a, b \in K$ and $w_1, w_2 \in W$. Put $v_1 := T^{-1}(w_1)$ and $v_2 := T^{-1}(w_2)$, so that $T(v_1) = w_1$ and $T(v_2) = w_2$. Since $T$ is linear, we have
\[
T(a v_1 + b v_2) = a T(v_1) + b T(v_2) = a w_1 + b w_2.
\]
Applying $T^{-1}$ to both sides gives
\[
T^{-1}(a w_1 + b w_2) = a v_1 + b v_2 = a \cdot T^{-1}(w_1) + b \cdot T^{-1}(w_2),
\]
so $T^{-1}$ is a linear map.
\end{proof}

\begin{summary}
\label{sum:bijective_equals_iso}
For linear maps, bijective $=$ isomorphism.
\end{summary}

% lem:15.b.5
\begin{lemma}[Compositions of Linear Maps are Linear]
\label{lem:composition_linear}
Let $T: V \to W$ and $S: W \to U$ be linear maps. Then the composition $S \circ T: V \to U$ is also linear.
\end{lemma}

\begin{proof}
Let $a, b \in K$ and $v_1, v_2 \in V$. Then
\begin{align*}
(S \circ T)(a v_1 + b v_2) &= S\big(T(a v_1 + b v_2)\big) = S\big(a T(v_1) + b T(v_2)\big) \\
&= a \cdot S(T(v_1)) + b \cdot S(T(v_2)) = a \cdot (S \circ T)(v_1) + b \cdot (S \circ T)(v_2). \qedhere
\end{align*}
\end{proof}

\begin{exercise}[Isomorphism is an Equivalence Relation]
\label{exc:iso_equivalence_relation}
Show that $\cong$, that is, being isomorphic, defines an equivalence relation on the set of vector spaces over $K$.
\end{exercise}

% def:15.b.6
\begin{definition}[Endomorphisms and Automorphisms]
\label{def:endomorphism}
A linear map $T: V \to V$, whose domain and target are the same space $V$, is sometimes called an \newterm{endomorphism} of $V$. We denote the set of all endomorphisms of $V$ by $\End(V)$, so that $\End(V) = \Hom(V, V)$. A linear map $T: V \to V$ which is an isomorphism is called an \newterm{automorphism} of $V$; in slogan form, endomorphism $+$ isomorphism $=$ automorphism.
\end{definition}

% def:15.b.7
\begin{definition}[Monomorphisms and Epimorphisms]
\label{def:monomorphism}
If a linear map $T$ is injective, one says that $T$ is a \newterm{monomorphism}; if $T$ is surjective, one says that $T$ is an \newterm{epimorphism}.
\end{definition}

\begin{remark}
We will not use the terminology of \cref{def:monomorphism} in this course, and will simply say injective and surjective instead.
\end{remark}

\subsection{The Rank Theorem}

% lem:15.b.8
\begin{lemma}[The Image is Spanned by the Images of a Basis]
\label{lem:image_spanned_by_basis}
Let $T: V \to W$ be a linear map, and let $v_1, \dots, v_n$ be a basis of $V$. Then
\[
\Img(T) = \Sp\big(T(v_1), \dots, T(v_n)\big).
\]
\end{lemma}

\begin{proof}
We first show that $\Img(T) \supseteq \Sp(T(v_1), \dots, T(v_n))$. Let $a_1, \dots, a_n \in K$. Then
\[
\sum_{i=1}^{n} a_i T(v_i) = T\left( \sum_{i=1}^{n} a_i v_i \right) \in \Img(T),
\]
which shows that $\Sp(T(v_1), \dots, T(v_n)) \subseteq \Img(T)$.

We show now that $\Img(T) \subseteq \Sp(T(v_1), \dots, T(v_n))$. Let $w \in \Img(T)$. By definition, there exists $v \in V$ such that $T(v) = w$. Since $v_1, \dots, v_n$ is a basis for $V$, there exist $a_1, \dots, a_n \in K$ such that $v = \sum_{i=1}^{n} a_i v_i$, and therefore,
\[
w = T(v) = T\left( \sum_{i=1}^{n} a_i v_i \right) = \sum_{i=1}^{n} a_i T(v_i).
\]
So $w \in \Sp(T(v_1), \dots, T(v_n))$, which proves the second inclusion.
\end{proof}

% Creator: Opus 5
The dimension of the image is the quantity the next theorem is about; it is given the name $\rank(T)$ at the end of this section, in \cref{def:rank_linear_map}.

% Extractor: Gemini 3.1 Pro
% thm:15.b.9
\begin{theorem}[Rank-Nullity Theorem]
\label{thm:rank_nullity}
Let $V$ be a finite-di\-men\-sional vector space, and let $T: V \to W$ be a linear map. Then, the following dimension formula holds:
\[
\dim V = \dim \ker(T) + \dim \Img(T).
\]
\end{theorem}

\begin{proof}
Being a subspace of the finite-di\-men\-sional space $V$, the kernel $\ker(T)$ is itself finite-di\-men\-sional by \cref{prop:subspace_dimension_constraints}, so it has a basis. Let $k := \dim \ker(T)$ and $n := \dim V$, and let $(u_1, \dots, u_k)$ be a basis for $\ker(T)$. By the basis extension property, item \textbf{(2)} of \cref{sum:finite_dim_properties}, we can extend this linearly independent set to a full basis for $V$, denoted by $\mathcal{B} = (u_1, \dots, u_k, v_1, \dots, v_{n-k})$. We claim that the family $\mathcal{C} = (T(v_1), \dots, T(v_{n-k}))$ forms a basis for $\Img(T)$.

\begin{enumerate}
    \item \textbf{Spanning:} Let $w \in \Img(T)$. By definition, there exists some $v \in V$ such that $w = T(v)$. We can express $v$ uniquely as a linear combination of the basis vectors in $\mathcal{B}$:
    \[
    v = \sum_{i=1}^k a_i u_i + \sum_{j=1}^{n-k} b_j v_j.
    \]
    Applying the linear map $T$, and recalling that $T(u_i) = 0_W$ for all $i$, we find:
    \begin{align*}
    w &= T\left( \sum_{i=1}^k a_i u_i + \sum_{j=1}^{n-k} b_j v_j \right) \\
    &= \sum_{i=1}^k a_i T(u_i) + \sum_{j=1}^{n-k} b_j T(v_j) \\
    &= 0_W + \sum_{j=1}^{n-k} b_j T(v_j).
    \end{align*}
    Thus, the vectors in $\mathcal{C}$ span $\Img(T)$.
    
    \item \textbf{Linear Independence:} Suppose there exist scalars $b_j \in K$ such that:
    \[
    \sum_{j=1}^{n-k} b_j T(v_j) = 0_W.
    \]
    By the linearity of $T$, this is equivalent to $T\left( \sum_{j=1}^{n-k} b_j v_j \right) = 0_W$, which implies that the vector $\sum_{j=1}^{n-k} b_j v_j$ lies in $\ker(T)$. Because $(u_1, \dots, u_k)$ is a basis for the kernel, there must exist scalars $a_i \in K$ such that:
    \[
    \sum_{j=1}^{n-k} b_j v_j = \sum_{i=1}^k a_i u_i,
    \]
    or, after moving everything to one side,
    \[
    \sum_{j=1}^{n-k} b_j v_j - \sum_{i=1}^k a_i u_i = 0_V.
    \]
    However, since $\mathcal{B}$ is a basis for $V$, all of its constituent vectors are linearly independent. This forces all coefficients to be zero, specifically $b_j = 0$ for all $1 \le j \le n-k$. 
\end{enumerate}

Therefore, $\mathcal{C}$ is linearly independent and spans the image, making it a basis for $\Img(T)$. Consequently, $\dim \Img(T) = n - k = \dim V - \dim \ker(T)$.
\end{proof}

% Extractor: Opus 5
\begin{remark}
The spanning step above can be shortened: applying \cref{lem:image_spanned_by_basis} to the basis $\mathcal{B}$ gives
\[
\Img(T) = \Sp\big(\underbrace{T(u_1)}_{= \, 0}, \dots, \underbrace{T(u_k)}_{= \, 0}, T(v_1), \dots, T(v_{n-k})\big) = \Sp\big(T(v_1), \dots, T(v_{n-k})\big),
\]
which is how the notes obtain it. Only the independence of $\mathcal{C}$ then remains to be checked.
\end{remark}

% Creator: Opus 5
The theorem also admits a second, more structural proof, which replaces the manipulation of coordinates by a single well-chosen decomposition of $V$. It uses the existence of a complement from \cref{prop:existence_complement} and the fact that isomorphisms carry bases to bases, \cref{thm:iso_preserves_bases} below; neither of these depends on the rank theorem, so no circularity arises.

\begin{proof}[Second proof of \Cref{thm:rank_nullity}]
By \cref{prop:existence_complement}, the subspace $\ker(T) \subseteq V$ admits a complement, that is, a subspace $U \subseteq V$ with
\[
\ker(T) + U = V \quad \text{and} \quad \ker(T) \cap U = \{0_V\}.
\]
We claim that the restriction $T|_U: U \to \Img(T)$ is an isomorphism.

It is injective: if $u \in U$ satisfies $T(u) = 0_W$, then $u \in \ker(T) \cap U = \{0_V\}$, so $u = 0_V$, and \cref{prop:injectivity_kernel} \textbf{(a)} applies. It is surjective onto $\Img(T)$: given $w \in \Img(T)$, write $w = T(v)$ for some $v \in V$ and decompose $v = x + u$ with $x \in \ker(T)$ and $u \in U$; then
\[
w = T(x + u) = T(x) + T(u) = 0_W + T(u) = T(u).
\]
Being a bijective linear map, $T|_U$ is an isomorphism by \cref{lem:bijective_is_iso}. By \cref{thm:iso_preserves_bases}, it carries a basis of $U$ to a basis of $\Img(T)$, and therefore, $\dim U = \dim \Img(T)$.

Finally, the dimension formula of \cref{prop:properties_of_sums} \textbf{(c)}, applied to the subspaces $\ker(T)$ and $U$, gives
\[
\dim V = \dim\big(\ker(T) + U\big) = \dim \ker(T) + \dim U - \underbrace{\dim\big(\ker(T) \cap U\big)}_{= \, 0} = \dim \ker(T) + \dim \Img(T). \qedhere
\]
\end{proof}

% cor:15.b.10
\begin{corollary}[Dimension Obstructions to Injectivity and Surjectivity]
\label{cor:equal_dim_iso}
% Extractor: Opus 5
Let $T: V \to W$ be a linear map, where $V$ and $W$ are finite-di\-men\-sional vector spaces. Then:
\begin{enumerate}[label=\textbf{(\alph*)}]
    \item If $\dim W < \dim V$, then $T$ is \textbf{not} injective.
    \item If $\dim W > \dim V$, then $T$ is \textbf{not} surjective.
    \item If $\dim W = \dim V$, then $T$ is bijective if and only if $T$ is injective, if and only if $T$ is surjective.
\end{enumerate}
\end{corollary}

\begin{proof}
\begin{enumerate}[label=\textbf{(\alph*)}]
    \item Since $\Img(T)$ is a subspace of $W$, we have $\dim \Img(T) \le \dim W < \dim V$ by \cref{prop:subspace_dimension_constraints}. By \cref{thm:rank_nullity},
    \[
    \dim \ker(T) = \dim V - \dim \Img(T) > 0,
    \]
    so $\ker(T) \neq \{0_V\}$, and $T$ is not injective by \cref{prop:injectivity_kernel} \textbf{(a)}.
    \item By \cref{thm:rank_nullity} and the assumption,
    \[
    \dim \Img(T) = \dim V - \dim \ker(T) \le \dim V < \dim W,
    \]
    so $\Img(T) \subsetneq W$, and $T$ is not surjective by \cref{prop:injectivity_kernel} \textbf{(b)}.
    \item Here the chain of equivalences reads
    \[
    T \text{ injective} \iff \ker(T) = \{0_V\} \iff \dim \ker(T) = 0 \iff \dim V = \dim \Img(T) \iff \dim W = \dim \Img(T),
    \]
    the third step by \cref{thm:rank_nullity} and the last by the assumption $\dim V = \dim W$. Since $\Img(T)$ is a subspace of $W$, the final equality holds if and only if $\Img(T) = W$, this being the equality case of \cref{prop:subspace_dimension_constraints}, that is, if and only if $T$ is surjective. As each of the two properties implies the other, either one already makes $T$ bijective. \qedhere
\end{enumerate}
\end{proof}

\begin{remark}
By \cref{lem:bijective_is_iso}, the word \qt{bijective} in statement \textbf{(c)} may equivalently be read as \qt{is an isomorphism}. So for maps between spaces of equal finite dimension, injectivity alone already forces $T$ to be an isomorphism.
\end{remark}

% cor:15.b.11
\begin{corollary}[Isomorphism is Detected by Dimension]
\label{cor:iso_iff_equal_dimension}
Two finite-di\-men\-sional vector spaces $V$ and $W$ over $K$ are isomorphic if and only if $\dim V = \dim W$.
\end{corollary}

\begin{proof}
Suppose first that $T: V \to W$ is an isomorphism. By \cref{def:isomorphism}, it admits a linear inverse and is in particular bijective, so $\ker(T) = \{0_V\}$ and $\Img(T) = W$ by \cref{prop:injectivity_kernel}. The rank theorem, \cref{thm:rank_nullity}, then gives
\[
\dim W = \dim \Img(T) = \dim V - \dim \ker(T) = \dim V.
\]

Conversely, assume that $\dim V = \dim W$, and denote this common dimension by $n$. Let $v_1, \dots, v_n$ be a basis for $V$ and $w_1, \dots, w_n$ a basis for $W$. By \cref{thm:existence_uniqueness_linear_map}, there is a linear map $T: V \to W$ with $T(v_i) = w_i$ for all $1 \le i \le n$. By \cref{lem:image_spanned_by_basis},
\[
\Img(T) = \Sp\big(T(v_1), \dots, T(v_n)\big) = \Sp(w_1, \dots, w_n) = W,
\]
so $T$ is surjective and $\dim \Img(T) = n$. By \cref{thm:rank_nullity}, $\dim \ker(T) = n - n = 0$, so $T$ is injective as well. Being a bijective linear map, $T$ is an isomorphism by \cref{lem:bijective_is_iso}, and hence $V \cong W$.
\end{proof}

\subsection{Isomorphisms and Vector Space Equivalence}

% Extractor: Gemini 3.1 Pro
Isomorphisms allow us to classify vector spaces and identify when two ostensibly different spaces are essentially the \qt{same} in terms of their linear algebraic structure.

% thm:15.b.12
\begin{theorem}[Isomorphisms Preserve Bases]
\label{thm:iso_preserves_bases}
Let $T: V \to W$ be an isomorphism between two vector spaces $V$ and $W$. Let $\mathcal{S} = \{v_i\}$ be a family of vectors in $V$. We define its image family as $T(\mathcal{S}) = \{T(v_i)\}_{v_i \in \mathcal{S}} \subseteq W$. Then, the following hold:
\begin{enumerate}
    \item $\mathcal{S}$ is linearly independent in $V$ if and only if $T(\mathcal{S})$ is linearly independent in $W$.
    \item $\mathcal{S}$ spans $V$ if and only if $T(\mathcal{S})$ spans $W$.
    \item $\mathcal{S}$ is a basis of $V$ if and only if $T(\mathcal{S})$ is a basis of $W$.
\end{enumerate}
\end{theorem}

% Creator: Opus 5
The notes state this theorem without proof. Since the second proof of
\cref{thm:rank_nullity} above appeals to it, we supply one; it is short, and it
uses nothing beyond the definition of an isomorphism.

\begin{proof}
Write $S := T^{-1}: W \to V$ for the inverse of $T$, which is linear because $T$ is an isomorphism. Each statement is an equivalence, and in each case the implication from right to left follows by applying the already-proved implication from left to right to the isomorphism $S$ and the family $T(\mathcal{S})$, since $S(T(\mathcal{S})) = \mathcal{S}$. It therefore suffices to prove the three implications from left to right.

\begin{enumerate}
    \item Let $\mathcal{S}$ be linearly independent, and suppose that a finite subfamily satisfies $\sum_{i} a_i T(v_i) = 0_W$. By linearity, $T\left( \sum_{i} a_i v_i \right) = 0_W$, and applying $S$ gives $\sum_{i} a_i v_i = 0_V$. The independence of $\mathcal{S}$ forces every $a_i$ to vanish, so $T(\mathcal{S})$ is linearly independent.
    \item Let $\mathcal{S}$ span $V$, and let $w \in W$. Then $S(w) \in V$ is a linear combination $S(w) = \sum_{i} a_i v_i$ of members of $\mathcal{S}$, and applying $T$ gives
    \[
    w = T(S(w)) = \sum_{i} a_i T(v_i),
    \]
    so $T(\mathcal{S})$ spans $W$.
    \item Immediate from the first two statements, since a basis is precisely a linearly independent spanning family. \qedhere
\end{enumerate}
\end{proof}

% Extractor: Gemini 3.1 Pro
\begin{remark}[The Significance of Isomorphism]
\label{rem:significance_of_iso}
If two vector spaces are isomorphic ($V \cong W$), they possess identical linear algebraic properties, such as having the exact same dimension. We can conceptually think of $V$ and $W$ as the exact same space, just represented in two different ways. 

However, it is crucial to note that identifying $V$ with $W$ requires the explicit choice of an isomorphism $T: V \to W$. Because there is usually no preferred (or \qt{canonical}) choice for this map, it is generally better practice not to forcefully identify different vector spaces simply because they are isomorphic. We will explore the nuances of canonical choices later in the course.
\end{remark}

\begin{exercise}[Isomorphisms Preserve Linear Independence, Spanning, and Bases]
\label{exc:isomorphisms_preserve_structures}
Let $V, W$ be vector spaces over a field $K$ and let $T \colon V \to W$ be an isomorphism. Prove directly that:
\begin{enumerate}
    \item $T$ maps every linearly independent subset of $V$ to a linearly independent subset of $W$;
    \item $T$ maps every spanning set of $V$ to a spanning set of $W$;
    \item $T$ maps every basis of $V$ to a basis of $W$.
\end{enumerate}
\exinfo{This exercise is Problem 6 of Problem Sheet 10, Linear Algebra I, Autumn Semester 2025.}
\end{exercise}

\begin{exercise}[Endomorphisms of One-Dimensional Spaces and Canonical Isomorphisms]
\label{exc:one_dim_endomorphism_canonical_iso}
Let $V$ be a one-dimensional vector space over a field $K$, and let $T \in \Hom_K(V, V)$.
\begin{enumerate}
    \item Show that there exists a unique scalar $\lambda \in K$ such that $T(v) = \lambda v$ for all $v \in V$.
    \item Explain why any isomorphism $V \to K$ depends essentially on a choice of a basis of $V$, whereas there is a natural (canonical) isomorphism $\Hom_K(V, V) \cong K$ that does not depend on any choice of basis.
\end{enumerate}
\exinfo{This exercise is Problem 1 of Problem Sheet 10, Linear Algebra I, Autumn Semester 2025.}
\end{exercise}

\begin{exercise}[Linear Functionals with the Same Kernel]
\label{exc:functionals_same_kernel_proportional}
Let $V$ be a vector space over a field $K$. Suppose that $T_1, T_2 \in \Hom_K(V, K)$ are two linear functionals such that $\ker(T_1) = \ker(T_2)$. Show that there exists a scalar $c \in K$ such that $T_1 = c T_2$.
\exhint{If the common kernel is $V$, both maps are zero. If not, consider a one-dimensional complement of the kernel and use \cref{exc:one_dim_endomorphism_canonical_iso}.}
\exinfo{This exercise is Problem 6 of Problem Sheet 11, Linear Algebra I, Autumn Semester 2025.}
\end{exercise}

\begin{exercise}[Projections and Complementary Subspaces]
\label{exc:projections_and_complements}
Let $V$ be a vector space. A linear endomorphism $P \colon V \to V$ is called a \newterm{projection} (or \newterm{idempotent}) if $P^2 := P \circ P = P$.
\begin{enumerate}
    \item Show that for every projection $P$, its image $\Img(P)$ is a linear complement of its kernel $\ker(P)$ in $V$, so that $V = \ker(P) \oplus \Img(P)$.
    \item Conversely, show that for any linear subspaces $W_1, W_2 \subseteq V$ such that $W_1$ is a complement of $W_2$ in $V$ ($V = W_1 \oplus W_2$), there exists a unique projection $P \colon V \to V$ with $\ker(P) = W_1$ and $\Img(P) = W_2$.
\end{enumerate}
\exinfo{This exercise is Problem 4 of Problem Sheet 10, Linear Algebra I, Autumn Semester 2025.}
\end{exercise}

\begin{exercise}[Invertibility from a Three-Map Composition]
\label{exc:invertibility_three_map_composition}
Suppose $V$ is a finite-dimensional vector space over $K$, and let $S, T, U \in \Hom(V, V)$ be endomorphisms satisfying
\[
S \circ T \circ U = \id_V.
\]
Show that $T$ is invertible and that its inverse is given by $T^{-1} = U \circ S$.
\exinfo{This exercise is Problem 2 of Problem Sheet 12, Linear Algebra I, Autumn Semester 2025.}
\end{exercise}

\begin{exercise}[Properties and Existence of Endomorphisms and Automorphisms]
\label{exc:endomorphism_properties_true_false}
Let $V$ be a finite-dimensional vector space over a field $K$. Prove or disprove each of the following statements:
\begin{enumerate}
    \item Let $V' \subseteq V$ be a linear subspace. Every automorphism $f \colon V' \to V'$ can be extended to an automorphism $\bar{f} \colon V \to V$ (i.e., $\bar{f}|_{V'} = f$).
    \item For every endomorphism $f \colon V \to V$, its image $\Img(f)$ is a linear complement of its kernel $\ker(f)$ in $V$.
    \item There does not exist any linear map $T \colon \mathbb{R}^5 \to \mathbb{R}^5$ such that $\rank(T) = \dim \ker(T)$.
\end{enumerate}
\exinfo{This exercise is Problem 4 of Problem Sheet 12, Linear Algebra I, Autumn Semester 2025.}
\end{exercise}

% def:15.b.13
\begin{definition}[Rank of a Linear Map]
\label{def:rank_linear_map}
Let $T \in \Hom(V, W)$ be a linear map such that $\dim \Img(T) < \infty$. We define the \newterm{rank} of $T$ as the dimension of its image:
\[
\rank(T) := \dim \Img(T).
\]
\end{definition}

% Extractor: Opus 5
\begin{exercise}[Rank of a Composition]
\label{exc:rank_of_composition}
Let $T: V \to W$ and $S: W \to U$ be linear maps, where $U$, $V$ and $W$ are finite-di\-men\-sional. Show that:
\begin{enumerate}[label=\textbf{(\alph*)}]
    \item $\rank(S \circ T) \le \min\{\rank(S), \rank(T)\}$;
    \item if $S$ is injective, then $\rank(S \circ T) = \rank(T)$;
    \item if $T$ is surjective, then $\rank(S \circ T) = \rank(S)$.
\end{enumerate}
\exinfo{This exercise is Problem 3 of Problem Sheet 10, Linear Algebra I, Autumn Semester 2025.}
\end{exercise}

% Creator: Opus 5
\subsection{Solutions to the Exercises}

\begin{exercisesolution}[Linearity Criterion (on \cpageref{exer:linearity_criterion})]
Suppose first that $T$ is linear. Then, for all $\alpha_1, \alpha_2 \in K$ and $v_1, v_2 \in V$,
\[
T(\alpha_1 v_1 + \alpha_2 v_2) = T(\alpha_1 v_1) + T(\alpha_2 v_2) = \alpha_1 T(v_1) + \alpha_2 T(v_2),
\]
using additivity for the first equality and homogeneity for the second.

Conversely, assume the displayed identity for all scalars and vectors. Taking $\alpha_1 = \alpha_2 = 1$ gives $T(v_1 + v_2) = T(v_1) + T(v_2)$, which is additivity. Taking $\alpha_2 = 0$ and $v_2 = 0_V$ gives
\[
T(\alpha_1 v_1) = \alpha_1 T(v_1) + 0 \cdot T(0_V) = \alpha_1 T(v_1),
\]
which is homogeneity. So $T$ satisfies both conditions of \cref{def:linear_map}.
\end{exercisesolution}

\begin{exercisesolution}[Images and Preimages of Subspaces (on \cpageref{exc:images_preimages_subspaces})]
\begin{enumerate}[label=\textbf{(\alph*)}]
    \item Since $0_V \in V'$ and $T(0_V) = 0_W$, we have $0_W \in T(V')$. Let $a, b \in K$ and $w_1, w_2 \in T(V')$, say $w_1 = T(v_1)$ and $w_2 = T(v_2)$ with $v_1, v_2 \in V'$. Then $a v_1 + b v_2 \in V'$, because $V'$ is a subspace, and
    \[
    a w_1 + b w_2 = a T(v_1) + b T(v_2) = T(a v_1 + b v_2) \in T(V').
    \]
    \item Since $T(0_V) = 0_W \in W'$, we have $0_V \in T^{-1}(W')$. Let $a, b \in K$ and $v_1, v_2 \in T^{-1}(W')$, so that $T(v_1), T(v_2) \in W'$. Then
    \[
    T(a v_1 + b v_2) = a T(v_1) + b T(v_2) \in W',
    \]
    because $W'$ is a subspace, and therefore $a v_1 + b v_2 \in T^{-1}(W')$.
    \item Both parts of \cref{prop:kernel_image} are special cases. Taking $V' := V$ in \textbf{(a)} gives that $\Img(T) = T(V)$ is a subspace of $W$, and taking $W' := \{0_W\}$ in \textbf{(b)} gives that $\ker(T) = T^{-1}(\{0_W\})$ is a subspace of $V$, exactly the two set-theoretic descriptions recorded in the remark on \cpageref{rem:set_theoretic_notation}.
\end{enumerate}
\end{exercisesolution}

\begin{exercisesolution}[Isomorphism is an Equivalence Relation (on \cpageref{exc:iso_equivalence_relation})]
\textbf{Reflexivity.} The identity $\id_V: V \to V$ is linear and satisfies $\id_V \circ \id_V = \id_V$, so it is an isomorphism, and $V \cong V$.

\textbf{Symmetry.} Let $V \cong W$, witnessed by an isomorphism $T: V \to W$ with linear inverse $S: W \to V$. Then $S$ is itself linear with linear inverse $T$, since $T \circ S = \id_W$ and $S \circ T = \id_V$. So $S$ is an isomorphism and $W \cong V$.

\textbf{Transitivity.} Let $T: V \to W$ and $T': W \to U$ be isomorphisms with inverses $S$ and $S'$. By \cref{lem:composition_linear}, both $T' \circ T: V \to U$ and $S \circ S': U \to V$ are linear, and
\[
(S \circ S') \circ (T' \circ T) = S \circ (S' \circ T') \circ T = S \circ \id_W \circ T = S \circ T = \id_V,
\]
and symmetrically $(T' \circ T) \circ (S \circ S') = \id_U$. So $T' \circ T$ is an isomorphism and $V \cong U$.
\end{exercisesolution}

\begin{exercisesolution}[Rank of a Composition (on \cpageref{exc:rank_of_composition})]
Throughout, we use the identity
\[
\Img(S \circ T) = \{S(T(v)) \mid v \in V\} = \{S(w) \mid w \in \Img(T)\} = S\big(\Img(T)\big),
\]
so that $\Img(S \circ T)$ is the image of the restricted linear map $S|_{\Img(T)}: \Img(T) \to U$.

\begin{enumerate}[label=\textbf{(\alph*)}]
    \item Since $\Img(T) \subseteq W$, we have $S(\Img(T)) \subseteq S(W) = \Img(S)$, and both are subspaces of $U$ by \cref{exc:images_preimages_subspaces} \textbf{(a)}, so \cref{prop:subspace_dimension_constraints} gives $\rank(S \circ T) \le \rank(S)$. Applying \cref{thm:rank_nullity} to $S|_{\Img(T)}$ gives
    \[
    \rank(S \circ T) = \dim \Img\big(S|_{\Img(T)}\big) = \dim \Img(T) - \dim \ker\big(S|_{\Img(T)}\big) \le \rank(T),
    \]
    and the two bounds together give $\rank(S \circ T) \le \min\{\rank(S), \rank(T)\}$.
    \item If $S$ is injective, then so is its restriction $S|_{\Img(T)}$, so $\ker(S|_{\Img(T)}) = \{0_W\}$ and the displayed computation becomes an equality: $\rank(S \circ T) = \dim \Img(T) = \rank(T)$.
    \item If $T$ is surjective, then $\Img(T) = W$, and therefore $\Img(S \circ T) = S(W) = \Img(S)$, so that $\rank(S \circ T) = \rank(S)$. \qedhere
\end{enumerate}
\end{exercisesolution}

\begin{ainote}
This was the most heavily damaged chapter the restoration met. Of the thirteen
numbered statements in the notes for \cref{sec:kernel_image_rank}, the transcript
kept five and hid the loss with a counter jump, a
\verb|\setcounter{theorem}{11}| annotated as aligning the numbering with the
notes. The missing eight are restored above in Prof. Biran's own order, so the
printed numbers now agree with the handwritten ones throughout.

The gap that mattered most was \cref{def:isomorphism}. The transcript carried a
definition of \emph{isomorphic spaces} at the very end of the chapter, while the
word \qt{isomorphism} was already in use well above it, and the old proof of
\cref{cor:equal_dim_iso} concluded from bijectivity that a map is an isomorphism,
which is exactly \cref{lem:bijective_is_iso}, the one result the transcript had
dropped. Both are back in the notes' positions.

Two additions go beyond the notes. Prof. Biran leaves
\cref{thm:iso_preserves_bases} unproved, so a proof is supplied; it then yields a
second proof of \cref{thm:rank_nullity}, which takes a complement $U$ of
$\ker(T)$ as in \cref{prop:existence_complement} and observes that $T|_U$ is an
isomorphism onto $\Img(T)$. Neither ingredient depends on the rank theorem, so
the second proof is not circular.
\end{ainote}
